\documentclass[twocolumn,a4paper,aps,amsmath,amssymb,floatfix,superscriptaddress]{revtex4-2}
\usepackage[utf8]{inputenc}
\usepackage{booktabs}
\usepackage{graphicx}
\usepackage{amsmath, amssymb, amsthm}
\usepackage{physics}
\usepackage[shortlabels]{enumitem}
\usepackage{cancel}
\usepackage{lastpage}
\usepackage{braket}
\usepackage{mathrsfs}
\usepackage{appendix}
\usepackage{listings}
\usepackage{xcolor}
\usepackage{graphicx}
\usepackage{dcolumn}
\usepackage{hyperref}
\usepackage{bm}
\usepackage[utf8]{inputenc}

\definecolor{codegreen}{rgb}{0,0.6,0}
\definecolor{codegray}{rgb}{0.5,0.5,0.5}
\definecolor{codepurple}{rgb}{0.58,0,0.82}
\definecolor{backcolour}{rgb}{0.95,0.95,0.92}

\newcommand{\projectauthors}{Bela F. Brunner (8002193), Adam Zich (0074187)}
\newcommand{\projectdate}{\today}
\newcommand{\maketitlepage}{
    \begin{titlepage}
        \centering
        
        \vspace*{0.2\textheight}

        {\Huge\bfseries Determining critical phase transition temperature in an Ising model simulation\par}
        \vspace{2cm}

        {\Large \projectauthors\par}
        \vspace{2cm}

        {\large \projectdate\par}

        \vfill
    \end{titlepage}
}

\newcommand{\todo}[1]{
    \textbf{TODO: $<$#1$>$}
}

\usepackage{fancyhdr}
\pagestyle{fancy}
\lhead{NS-204B - Statystische fysica theorie en experiment}
\rhead{\projectauthors}
\fancyfoot{}
\fancyfoot[R]{Page \thepage /\pageref{LastPage}}

\newcommand{\R}{\mathbb{R}}
\newcommand{\p}{\mathbb{P}}
\newcommand{\N}{\mathbb{N}}
\newcommand{\Z}{\mathbb{Z}}
\newcommand{\Q}{\mathbb{Q}}
\newcommand{\I}{\mathbb{I}}
\newcommand{\E}{\mathbb{E}}

\newcommand{\mbeq}{\overset{!}{=}}
\newcommand{\ko}{\frac{1}{4\pi \epsilon_0}}

\lstset{frame=tb,
  language=java,
  aboveskip=3mm,
  belowskip=3mm,
  showstringspaces=false,
  columns=flexible,
  basicstyle={\small\ttfamily},
  numbers=none,
  numberstyle=\tiny\color{gray},
  keywordstyle=\color{blue},
  commentstyle=\color{dkgreen},
  stringstyle=\color{mauve},
  breaklines=true,
  breakatwhitespace=true,
  tabsize=3
}

\lstdefinestyle{mystyle}{
    backgroundcolor=\color{backcolour},   
    commentstyle=\color{codegreen},
    keywordstyle=\color{blue},
    numberstyle=\tiny\color{codegray},
    stringstyle=\color{codegreen},
    basicstyle={\small\ttfamily},,
    breakatwhitespace=false,         
    breaklines=true,                 
    captionpos=b,                    
    keepspaces=true,                 
    numbers=left,                    
    numbersep=5pt,                  
    showspaces=false,                
    showstringspaces=false,
    showtabs=true,                  
    tabsize=1
}

\lstset{style=mystyle}

\begin{document}

%% Title and authors using built-in LaTeX options. Change according to own preferences
\title{Determining critical phase transition temperature in an Ising model simulation}
        
\author{Bela F. Brunner}
%Lines break automatically or can be forced with \\
\affiliation{Department of Physics, \\
Utrecht University}

\author{Adam Zich}
\affiliation{Department of Physics, \\
Utrecht University}

\date{\today}
% It is always \today, today,
% but any date may be explicitly specified

\begin{abstract}

Room for abstract.

\end{abstract}

\maketitle

\section{Introduction}\label{intro}
    \subsection{The Ising model}
        The Ising model named after the German physicist Erns Ising is a model in statistical mechanics. The model consists of discrete magnetic spin dipole moments that take values $\pm 1$ ("up" and "down" respectively) arranged on a lattice. For the purposes of our simulation we used a square lattice. The spins interact with neighbouring spins. In our implementation we consider the von Neumann neighbourgood of only 4 neighbouring spins in the $\pm x$ and $\pm y$ directions. The model exhibits phase transitions that can be associated with the gas-liquid transition or spontaneous magnetisation.
    \subsection{Aim of study}
        In this article we look at determining at what temperature does the Ising model undergo a phase transition. It is know from theory that under a certain critical temperature the model stabilises into completely up or down state. Conversely, above this critical, on average, the model remains a random mixture of up and down spins.
    
\section{Theory}\label{theory}

Theory.
% When the Theory is not the new aspect of the research (as is mostly the case in course experiments or experimental research in general), keep it short. Describing a starting point (e.g. a differential equation or a model), a final result (e.g. a relation under test, or a dependance) and a reference might be sufficient. 

\section{Methodology}
    \subsection{Metropolis algorithm}
        \todo{Explain metropolis algorithm}
    \subsection{Determining parameters for final simulation}
        \todo{Explain how we arrived at simulation parameters}
    

% Note that the method section about a simulation differs from a setup and method section of a physical experiment.

% The goal of the method section is to decribe the procedure of the experiment in such a way that a reader could reproduce the experiment (either to verify your results, or to build future research on your findings). This goal remains the same.

% Instead of a diagram of the measurement setup you could include a schematic of the interactions or steps in your model.

% Do not copy your code into the method section, instead describe the functions and actions of your code, and describe which packages or built-in functions you use.

\section{Results}\label{results}
    \todo{Give results}\\
    \todo{Give figures of results}
    
\section{Comparison to theory}
    \subsection{Theoretical solution for the critical phase transition temperature}
        \todo{Give derivation of critical temperature in the Ising model}
    \subsection{Comparison of theory and our results}
        \todo{Compare our results with theoreticla values}

\section{Conclusions}\label{conclusions}
Conclusion.
% Summarize the final results and answer the research question. Reflect on the strength and impact of your conclusion.
% In principle no new information or arguments are presented in the conclusion.


\bibliographystyle{unsrt}
\addcontentsline{toc}{section}{Bibliography}
%\bibliography{bibliography}            % *.bib Files

%Add bbl file
\begin{thebibliography}{10}

\bibitem{chen_JAP_2003}
F.~Chen, M.C. Cheung, P.M. Sweeney, W.D. Kirkey, M.~Furis, and A.N. Cartwright.
%\newblock Ultrafast differential transmission spectroscopy of excitonic
%  transitions in {I}n{G}a{N}/{G}a{N} multiple quantum wells.
\newblock J. Appl. Phys. \textbf{93}, 4933 (2003).

% You can add and cite as many new bibitems as you want; make sure to give them a unique and reccognizable name.

\end{thebibliography}


\end{document}
%
% ****** End of file ******
